\documentclass[11pt]{article}

\usepackage[margin=1in]{geometry}
\usepackage{amsmath,amssymb}
\usepackage{braket}
\usepackage[colorlinks=true, linkcolor=blue, citecolor=blue, urlcolor=blue]{hyperref}
\usepackage{graphicx}

\title{\textbf{Quench Dynamics of the 1D XXZ Model}\\
\large Finite-Size Parity Structure and Symmetry-Broken States}

\author{Ho Jang}

\date{\today}

\begin{document}

\maketitle

\section{Weak random field and finite-size ED}

Consider the XXZ chain with a weak random field only in the initial Hamiltonian,
\begin{align}
    H_i = H_{\rm XXZ}(\Delta_i) + V,
    \qquad
    V=-\sum_j h_j S_j^z ,
    \qquad
    |h_j|\ll 1 .
\end{align}
After preparing the ground state of $H_i$, the system evolves with the clean final
Hamiltonian
\begin{align}
    H_f = H_{\rm XXZ}(\Delta_f).
\end{align}
The clean XXZ Hamiltonian has the global spin-flip symmetry
\begin{align}
    P H_{\rm XXZ} P^\dagger = H_{\rm XXZ},
    \qquad
    P S_j^z P^\dagger = -S_j^z .
\end{align}
Therefore finite-size ED eigenstates can be chosen as parity eigenstates, not particular symmetry broken state.

\section{Initial state}

Let $\ket{GS_i^+}$ be the even-parity ground state of the clean initial Hamiltonian.
Since the random field is parity odd,
\begin{align}
    P V P^\dagger = -V ,
\end{align}
it mixes the even ground state with odd-parity states which will be represented as $\ket{k^-}$. To first order when the state is non-degenerate(well-defined denominator), 
\begin{align}
    \ket{\psi(0)}
    \simeq
    \ket{GS_i^+}
    +
    \sum_{k^-}
    \frac{\bra{k^-}V\ket{GS_i^+}}
    {E_{GS_i}^+-E_k^-}
    \ket{k^-}.
    \label{eq:initial_state}
\end{align}
If $V \ll 1$ with weak random field magnitude, the odd-sector weight is small. In this case, the initial state is not a fully symmetry-broken N\'eel state. It is mostly
the even-parity state plus a small odd-parity correction. For $\Delta_i\leq 1$, the initial state is not in the ordered Ising regime.

\section{After the quench}

The final Hamiltonian also preserves parity. Expanding the initial state in the eigenbasis
of $H_f$,
\begin{align}
    \ket{\psi(0)}
    =
    \sum_n c_n^+ \ket{n^+}
    +
    \sum_m c_m^- \ket{m^-},
\end{align}
with
\begin{align}
    c_n^+ = O(1),
    \qquad
    c_m^- = O(h).
\end{align}
The time-evolved state is
\begin{align}
    \ket{\psi(t)}
    =
    \sum_n c_n^+ e^{-iE_n^+t}\ket{n^+}
    +
    \sum_m c_m^- e^{-iE_m^-t}\ket{m^-}.
\end{align}

Since $S_i^z$ is parity odd, same-parity matrix elements vanish:
\begin{align}
    \bra{n^+}S_i^z\ket{n'^+}=0,
    \qquad
    \bra{m^-}S_i^z\ket{m'^-}=0.
\end{align}
Therefore,
\begin{align}
    \langle S_i^z(t)\rangle
    =
    2\,{\rm Re}
    \sum_{n,m}
    (c_n^+)^*c_m^-
    \bra{n^+}S_i^z\ket{m^-}
    e^{i(E_n^+-E_m^-)t}.
    \label{eq:local_sz}
\end{align}

This equation explains the ED result. The matrix element
\begin{align}
    \bra{n^+}S_i^z\ket{m^-}
\end{align}
can contain the N\'eel spatial structure, but the amplitude is controlled by
\begin{align}
    (c_n^+)^*c_m^- \sim O(h).
\end{align}
Thus, for a perturbative initial random field, the magnetization scale along it, which is shown in Fig.~\ref{fig:ed-xxz}.
\begin{align}
    \boxed{
    \langle S_i^z(t)\rangle \sim O(h)
    } .
\end{align}

\begin{figure}
    \centering
    \includegraphics[width=0.99\linewidth]{benchmark-two-site-tdvp/ED-xxz-quench.pdf}
    \caption{XXZ quench with $\Delta_i=0.5$ with random magnetic field to $\Delta_f = 2.0$ without external magnetic field. The strength of Neel order depends on the initial random disorder under perturbative consideration.}
    \label{fig:ed-xxz}
\end{figure}

\section{TDVP result}
Here I compare ED and TDVP using the same quench protocol and the same random-field seed. 
For finite initial random-field strength, the agreement is good, as shown in Fig.~\ref{fig:tdvp-xxz}.

\begin{figure}[h!]
    \centering
    \includegraphics[width=0.99\linewidth]{benchmark-two-site-tdvp/TDVP-xxz-quench.pdf}
    \caption{
    TDVP result under the same condition as ED, using the same random-field seed. 
    The finite-field cases show consistent magnetization scales between ED and TDVP.
    }
    \label{fig:tdvp-xxz}
\end{figure}

This suggests that both ED and TDVP capture the same parity-controlled response in this regime.




\section{Symmetry-broken state}

In nature, one of the two N\'eel branches is selected in the ordered regime. \textbf{I think large $\Delta_f$, $N$ can break symmetry.}

%In small-size ED, however, the clean Hamiltonian still preserves the global spin-flip
%symmetry, so the exact low-energy states are parity eigenstates rather than explicitly
%symmetry-broken states.

Let
\begin{align}
    \ket{A}
    =
    \ket{\uparrow\downarrow\uparrow\downarrow\cdots},
    \qquad
    \ket{B}
    =
    \ket{\downarrow\uparrow\downarrow\uparrow\cdots}
\end{align}
denote the two N\'eel-like branches. For a finite clean system under parity symmetry, the corresponding low-energy
states are approximately
\begin{align}
    \ket{\rm even}
    =
    \frac{1}{\sqrt{2}}
    \left(
    \ket{A}
    +
    \ket{B}
    \right),
    \qquad
    \ket{\rm odd}
    =
    \frac{1}{\sqrt{2}}
    \left(
    \ket{A}
    -
    \ket{B}
    \right).
\end{align}
If $\ket{A}$ and $\ket{B}$ were completely disconnected, these two states would be exactly
degenerate. In a finite quantum system, however, quantum fluctuations weakly connect the
two branches. Restricting to the subspace spanned by $\ket{A}$ and $\ket{B}$, the effective
Hamiltonian can be written as
\begin{align}
    H_{\rm eff}
    =
    \begin{pmatrix}
        E_0 & t \\
        t & E_0
    \end{pmatrix},
    \label{eq:two_state_heff}
\end{align}
where $t$ is the finite-size tunneling amplitude between the two N\'eel branches. 
Diagonalizing this two-state Hamiltonian gives the even--odd energy splitting
\begin{align}
    |\Delta_{\rm eo}(N)|
    =
    |E_{\rm odd}-E_{\rm even}|
    =
    2|t|.
\end{align}
The transverse part of hamiltonian acts only locally,
\begin{align}
    H_\perp
    =
    \frac{J}{2}
    \sum_i
    \left(
        S_i^+S_{i+1}^-
        +
        S_i^-S_{i+1}^+
    \right).
\end{align} 
One application of $H_\perp$ cannot convert $\ket{A}$ into
$\ket{B}$. Instead, it creates a local defect, or equivalently a pair of domain
walls. Further local flips move these domain walls through the chain. Only after
a sequence of $O(N)$ local virtual processes can the system connect the two
N\'eel branches.

Schematically,
\begin{align}
    \ket{A}
    \rightarrow
    \ket{\text{domain-wall state}}
    \rightarrow
    \ket{\text{domain wall moved}}
    \rightarrow
    \cdots
    \rightarrow
    \ket{B}.
\end{align}
Each intermediate domain-wall configuration costs an Ising energy of order
\begin{align}
    \Delta_{\rm DW} \sim J\Delta,
\end{align}
while each local spin-exchange matrix element is of order $J$. Therefore each
virtual step contributes a small factor
\begin{align}
    \frac{J}{\Delta_{\rm DW}}
    \sim
    \frac{1}{\Delta}.
\end{align}
Since the number of required local steps grows with system size, the tunneling
amplitude behaves schematically as
\begin{align}
    t
    \sim
    J
    \left(
        \frac{1}{\Delta}
    \right)^{cN},
\end{align}
where $c$ is an order-one constant depending on boundary conditions and the
precise tunneling path. Thus
\begin{align}
    t
    \sim
    e^{-N/\xi},
\end{align}
with a short correlation length $\xi$ in the deep ordered regime. Therefore the even--odd splitting becomes exponentially small with system size. 

This also clarifies when the previous non-degenerate perturbation theory breaks down.
The random field couples the even and odd states through
\begin{align}
    v
    =
    \bra{\rm even}
    V
    \ket{\rm odd}.
\end{align}
If
\begin{align}
    |v| \ll |\Delta_{\rm eo}(N)|,
\end{align}
then the random field only weakly mixes the odd sector, and the local magnetization remains
perturbative. If instead
\begin{align}
    |v| \gtrsim |\Delta_{\rm eo}(N)|,
\end{align}
then the random field is nonperturbative within the lowest even--odd subspace. In that
case, the state becomes closer to one of the two symmetry-broken branches, $\ket{A}$ or
$\ket{B}$.

This is why \textbf{increasing $N,\Delta$} can help TDVP/DMRG show a symmetry-broken profile. Under the existence of initial random disorder, as $N,\Delta$
increases, the even--odd splitting becomes very small, so a tiny physical or numerical bias
can select one low-entanglement N\'eel branch. However, larger N, we can't explore with ED. So we need TDVP.
\end{document}